%========================%
%     Bahasa Inggris     %
%========================%

% --- Skema Pembelajaran ---
\newcommand{\TutorialOnline}{Online Tutorial}
\newcommand{\TutorialWebinar}{Webinar Tutorial}
\newcommand{\TutorialTatapMuka}{Face-to-Face Tutorial}
%
% --- Data ---
%
\newcommand{\LembarJawaban}{Answer Sheet}
\newcommand{\Tugas}{Assignment}
\newcommand{\Sesi}{Session}
\newcommand{\Diskusi}{Discussion}
\newcommand{\DiskusiKe}{Discussion}
%
\newcommand{\DisusunOleh}{Prepared by}
\newcommand{\Nama}{Name}
\newcommand{\NomorIndukM}{Student ID}
\newcommand{\ProgramStudi}{Study Program}
\newcommand{\Fakultas}{Faculty}
\newcommand{\TutorPengampu}{Tutor}
\newcommand{\MataKuliah}{Subject/Course}
\newcommand{\KelasOnlineKe}{Class}
\newcommand{\UPBJJ}{Regional Office}
%
% --- Judul Biasa ---
%
\newcommand{\Soal}{Questions}
\newcommand{\Jawaban}{Answers}
%
% --- Judul Ilmiah ---
%
\newcommand{\KataPengantar}{Preface}
% --- Bab/Bagian 1
\newcommand{\Pendahuluan}{Introduction}
\newcommand{\LatarBelakang}{Background}
\newcommand{\LatarBelakangPenelitian}{Background of the Study}
\newcommand{\RumusanMasalah}{Problem Statement}
\newcommand{\RumusanMasalahPoin}{Research Questions}
\newcommand{\Tujuan}{Objectives}
\newcommand{\TujuanPenelitian}{Research Objectives}
% --- Bab/Bagian 2
\newcommand{\KajianTeori}{Theoretical Framework}
\newcommand{\KajianPustaka}{Literature Review}
\newcommand{\Metode}{Methods}
\newcommand{\MetodePenelitian}{Research Methods}
\newcommand{\Metodologi}{Methodology}
\newcommand{\MetodologiPenelitian}{Research Methodology}
% --- Bab/Bagian 3
\newcommand{\Pembahasan}{Discussion}
\newcommand{\HasilDanPembahasan}{Results and Discussion}
\newcommand{\TemuanDanPembahasan}{Findings and Discussion}
% --- Bab/Bagian 4
\newcommand{\Penutup}{Closure}
\newcommand{\Kesimpulan}{Conclusions}
\newcommand{\Saran}{Suggestions}
\newcommand{\Rekomendasi}{Recommendations}
% --- Lampiran
\newcommand{\LampiranSatu}{Appendix}
\newcommand{\LampiranBanyak}{Appendices}
%
% --- Teorema Matematika ----
%
\newcommand{\Teorema}{Theorem}
\newcommand{\Lema}{Lemma}
\newcommand{\Definisi}{Definition}
\newcommand{\Contoh}{Example}
%
% --- Nama Lampiran ---
%
%     Penamaan untuk gambar, tabel, matematika, dan kode yang ada
%     di keterangan lampiran maupun dirujuk di dalam kalimat; 
%     contohnya
%     lihat Gambar 2, lihat Tabel 3, Persamaan 4 menunjukkan, 
%     lihat Kode 5.
%
%     Bawaan Tanpa Disetel: 
%     Figure, Table, Equation, Listing
%
%     Bahasa Inggris: 
%     Figure, Table, Equation, Code
%
\ifundef{\captionsamerican}{}{
    \addto\captionsamerican{
        \renewcommand{\lstlistingname}{Code}
    }
}
\ifundef{\captionsbritish}{}{
    \addto\captionsbritish{
        \renewcommand{\lstlistingname}{Code}
    }
}



%==========================%
%     Bahasa Indonesia     %
%==========================%

\ifundef{\captionsindonesian}{}{
    % --- Skema Pembelajaran ---
    \renewcommand{\TutorialOnline}{Tutorial \textit{Online}}
    \renewcommand{\TutorialWebinar}{Tutorial \textit{Webinar}}
    \renewcommand{\TutorialTatapMuka}{Tutorial Tatap Muka}
    %
    % --- Data ---
    %
    \renewcommand{\LembarJawaban}{Lembar Jawaban}
    \renewcommand{\Tugas}{Tugas}
    \renewcommand{\Sesi}{Sesi}
    \renewcommand{\Diskusi}{Diskusi}
    \renewcommand{\DiskusiKe}{Diskusi ke-}
    %
    \renewcommand{\DisusunOleh}{Disusun Oleh}
    \renewcommand{\Nama}{Nama}
    \renewcommand{\NomorIndukM}{NIM}
    \renewcommand{\ProgramStudi}{Program Studi}
    \renewcommand{\Fakultas}{Fakultas}
    \renewcommand{\TutorPengampu}{Tutor Pengampu}
    \renewcommand{\MataKuliah}{Mata Kuliah}
    \renewcommand{\KelasOnlineKe}{Kelas \textit{Online} ke-}
    \renewcommand{\UPBJJ}{UPBJJ}
    %
    % --- Judul Biasa ---
    %
    \renewcommand{\Soal}{Soal}
    \renewcommand{\Jawaban}{Jawaban}
    %
    % --- Judul Ilmiah ---
    %
    \renewcommand{\KataPengantar}{Kata Pengantar}
    % --- Bab/Bagian 1
    \renewcommand{\Pendahuluan}{Pendahuluan}
    \renewcommand{\LatarBelakang}{Latar Belakang}
    \renewcommand{\LatarBelakangPenelitian}{Latar Belakang Penelitian}
    \renewcommand{\RumusanMasalah}{Rumusan Masalah}
    \renewcommand{\RumusanMasalahPoin}{Rumusan Masalah}
    \renewcommand{\Tujuan}{Tujuan}
    \renewcommand{\TujuanPenelitian}{Tujuan Penelitian}
    % --- Bab/Bagian 2
    \renewcommand{\KajianTeori}{Kajian Teori}
    \renewcommand{\KajianPustaka}{Kajian Pustaka}
    \renewcommand{\Metode}{Metode}
    \renewcommand{\MetodePenelitian}{Metode Penelitian}
    \renewcommand{\Metodologi}{Metodologi}
    \renewcommand{\MetodologiPenelitian}{Metodologi Penelitian}
    % --- Bab/Bagian 3
    \renewcommand{\Pembahasan}{Pembahasan}
    \renewcommand{\HasilDanPembahasan}{Hasil dan Pembahasan}
    \renewcommand{\TemuanDanPembahasan}{Hasil Temuan dan Pembahasan}
    % --- Bab/Bagian 4
    \renewcommand{\Penutup}{Penutup}
    \renewcommand{\Kesimpulan}{Kesimpulan}
    \renewcommand{\Saran}{Saran}
    \renewcommand{\Rekomendasi}{Rekomendasi}
    % --- Lampiran
    \renewcommand{\LampiranSatu}{Lampiran}
    \renewcommand{\LampiranBanyak}{Lampiran}
    %
    % --- Teorema Matematika ----
    %
    \renewcommand{\Teorema}{Teorema}
    \renewcommand{\Lema}{Lema}
    \renewcommand{\Definisi}{Definisi}
    \renewcommand{\Contoh}{Contoh}
    %
    % --- Nama Bulan dalam Bahasa Indonesia ---
    %
    %     Pengubahan nama bulan dalam Babel Bahasa Indonesia yang masih
    %     menggunakan pelafalan lokal, seperti Pebuari dan Nopember.
    %
    \renewcommand{\dateindonesian}{%
        \def\today{\number\day~\ifcase\month\or
            Januari\or Februari\or Maret\or April\or Mei\or Juni\or
            Juli\or Agustus\or September\or Oktober\or November\or Desember\fi
            \space\number\year}%
    }
    
    \addto\captionsindonesian{
        %
        % --- Judul Abstrak ---
        %
        %     Abstrak dalam Babel Bahasa Indonesia menggunakan
        %     "Ringkasan" daripada "Abstrak"
        %
        %     Anda dapat mengembalikannya menjadi "Ringkasan" jika mau
        %
        \renewcommand{\abstractname}{Abstrak}
        %
        % --- Judul Bibliografi/Daftar Pustaka/Referensi ---
        %
        % Bibliography ---
        % 
        %     "Bibliography" digunakan untuk document class yang 
        %     menggunakan Bab. "Reference" digunakan untuk document class
        %     tanpa Bab.
        %
        %     Judul bawaannya adalah "Bibliography"
        %     Judul bawaan dalam Bahasa Indonesia adalah "Bibliografi"
        %     Judul untuk Bahasa Indonesia pada template ini adalah
        %     "Referensi"
        %
        %     Anda juga bisa menggantinya menjadi "Daftar Pustaka" atau
        %     judul lain
        %
        \renewcommand{\bibname}{Referensi} % Bibliography
        \renewcommand{\refname}{Referensi} % Reference
        %
        % --- Nama Level Bagian untuk \autoref ---
        %
        %     Penamaan untuk level Chapter hingga Subparagraph saat 
        %     dirujuk
        %     di dalam kalimat, contohnya lihat Bagian 2.5.
        %
        %     [NOTE] Rujukan level bagian hanya efektif untuk penomoran 
        %            heading Multilevel (bernomor 1, 1.1, 1.1.1, 1.1.1.1)
        %
        %     Bawaan Tanpa Disetel:
        %     Chapter > Section > Subsection > Subsubsection > Paragraph %     > Subparagraph
        %
        %     Bahasa Indonesia (Terjemahan Paksa):
        %     Bab > Seksi > Subseksi > Subsubseksi > Paragraf >
        %     Subparagraf
        %
        %     Bahasa Indonesia (Terjemahan Paksa 2):
        %     Bab > Bagian > Subbagian > Subsubbagian > Paragraf > 
        %     Subparagraf
        %
        %     Bahasa Indonesia (Penyesuaian 1):
        %     Bab > Subbab > Bagian > Bagian > Bagian > Paragraf > 
        %     Subparagraf
        %
        %     Bahasa Indonesia (Penyesuaian 2) [Digunakan]:
        %     Bab > Bagian > Bagian > Bagian > Bagian > Paragraf > 
        %     Subparagraf
        %
        \renewcommand{\chapterautorefname}{Bab}
        \renewcommand{\sectionautorefname}{Bagian}
        \renewcommand{\subsectionautorefname}{Bagian}
        \renewcommand{\subsubsectionautorefname}{Bagian}
        \renewcommand{\paragraphautorefname}{Paragraf}
        \renewcommand{\subparagraphautorefname}{Subparagraf}
        %
        % --- Nama Lampiran ---
        %
        %     Penamaan untuk gambar, tabel, matematika, dan kode yang ada
        %     di keterangan lampiran maupun dirujuk di dalam kalimat; 
        %     contohnya lihat Gambar 2, lihat Tabel 3, Persamaan 4 
        %     menunjukkan, lihat Kode 5.
        %
        %     Bahasa Indonesia:
        %     Gambar, Tabel, Persamaan, Kode
        %
        \renewcommand{\figureautorefname}{Gambar}
        \renewcommand{\tableautorefname}{Tabel}
        \renewcommand{\equationautorefname}{Persamaan}
        \renewcommand{\lstlistingname}{Kode}
        %
        % --- Nama Lanjutan Tabel Tabularray di Halaman Berikutnya ---
        %
        %     Penamaan untuk "Continued to next page..." dan "(Continued)"
        %     pada Tabel Tabularray longtblr
        %
        %     Bawaan Tanpa Disetel: Continued to next page --- (Continued)
        %
        %     Bahasa Indonesia: Lanjutan di halaman berikutnya --- (Lanjutan)
        %
        \DeclareTblrTemplate{contfoot-text}{normal}{Lanjutan di halaman berikutnya}
        \SetTblrTemplate{contfoot-text}{normal}
        \DeclareTblrTemplate{conthead-text}{normal}{(Lanjutan)}
        \SetTblrTemplate{conthead-text}{normal}
    }
}